\documentclass[12pt]{article}
\usepackage[T1]{fontenc}
\usepackage[utf8]{inputenc}
\usepackage{amsmath,amssymb,mathrsfs}
\usepackage{pst-all,pst-func,pst-3dplot,pst-eucl}
\usepackage{multido,color}
\newcommand{\R}{\mathbb{R}}
\newcommand{\N}{\mathbb{N}}
\newcommand{\D}{\mathbb{D}}
\newcommand{\Z}{\mathbb{Z}}
\newcommand{\Q}{\mathbb{Q}}
\newcommand{\C}{\mathbb{C}}
\newcommand{\vect}[1]{\overrightarrow{\,\mathstrut#1\,}}
\newcommand{\barre}[1]{\overline{\,\mathstrut#1\,}}
\newcommand{\pg}{\geqslant}
\newcommand{\pp}{\leqslant}
\newcommand{\e}{\,\text{e}\,}
\renewcommand{\d}{\,\text{d}}
\newcommand{\Cg}{\texttt{]}}
\newcommand{\Cd}{\texttt{[}}
\def\Oij{$\left(\text{O}\,;\,\vec{\imath},\,\vec{\jmath}\right)$}
\def\Oijk{$\left(\text{O}\,;\,\vec{\imath},\,\vec{\jmath},\,\vec{k}\right)$}
\pagestyle{empty}
\begin{document}
\begin{center}
%		\begin{tikzpicture}%compiler avec pgfplots
%			\begin{axis}[/pgf/number format/.cd, use comma, x={45mm}, y={45mm}, xmin=-0.2,  xmax=2.2,    ymin = -0.2,     ymax=1.2, xtick ={1,2},    ytick={0.5,1},
%				tick label style={font=\footnotesize}, grid=none, axis lines =center]
%				\addplot [line width=1.2pt,color=red,smooth,samples=100,domain= 0.6:2.2 ]{ln(x)};
%				\node[below left] at (axis cs: 0,0) {{\footnotesize 0}};
%				\node[red] at(axis cs: 1.9,0.58) {$ \mathscr{C}$};
%				\addplot [line width=1pt,color=red,blue,samples=10,domain= 0.6:2.2 ]{x-1};
%				\node [blue] at(axis cs: 1.9,0.8) {T};
%				\fill (axis cs: 2,0) circle (1pt) node[above] {A$_0$};
%				\fill (axis cs: 1.414,0) circle (1pt) node[above] {A$_1$};
%				\fill (axis cs: 1.189,0) circle (1pt) node[above] {A$_2$};
%			\end{axis}
%		\end{tikzpicture}
\psset{xunit=4.5cm, yunit=4.5cm,arrowsize=2pt 3}
\begin{pspicture*}(-0.2,-0.2)(2.2,1.2)
\psaxes[linewidth=0.8pt,labels=none,Dx=1,Dy=0.5]{->}(0,0)(-0.2,-0.2)(2.2,1.2)
% Graduations manuelles
\psline[linewidth=0.5pt](1,-0.05)(1,0.05)
\rput[t](1,-0.1){\footnotesize 1}
\psline[linewidth=0.5pt](2,-0.05)(2,0.05)
\rput[t](2,-0.1){\footnotesize 2}
\psline[linewidth=0.5pt](-0.05,0.5)(0.05,0.5)
\rput[r](-0.1,0.5){\footnotesize 0{,}5}
\psline[linewidth=0.5pt](-0.05,1)(0.05,1)
\rput[r](-0.1,1){\footnotesize 1}
\rput[tr](-0.05,-0.05){\footnotesize 0}
% Courbe logarithme en rouge
\psplot[linewidth=1.25pt,linecolor=red,plotpoints=2000]{0.6}{2.2}{x ln}
\rput[l](1.9,0.58){\red $\mathscr{C}$}
\psplot[linewidth=1pt, linecolor=blue, plotpoints=10]{0.6}{2.2}{x 1 sub}
\rput[l](1.9,0.8){\blue T}
% Points
\psdot[dotsize=2pt, linecolor=black](2,0)
\rput[b](2,0.05){A$_0$}
\psdot[dotsize=2pt, linecolor=black](1.414,0)
\rput[b](1.414,0.05){A$_1$}
\psdot[dotsize=2pt, linecolor=black](1.189,0)
\rput[b](1.189,0.05){A$_2$}
\end{pspicture*}
\end{center}
\end{document}
